\documentclass[11pt,a4paper]{article}

% ---------- Packages ----------
\usepackage[margin=1in, top=1.2in, bottom=1.2in]{geometry}
\usepackage{tabularx}
\usepackage{array}
\usepackage{booktabs}
\usepackage{caption}
\usepackage[T1]{fontenc}
\usepackage{helvet}
\renewcommand{\familydefault}{\sfdefault}
\usepackage{microtype}
\usepackage{titling}
\usepackage{fancyhdr}
\usepackage{parskip}

% ---------- Header/Footer ----------
\pagestyle{fancy}
\fancyhf{}
\renewcommand{\headrulewidth}{0.5pt}
\fancyhead[L]{\small\textbf{Related Work Survey}}
\fancyhead[R]{\small WebRTC, P2P \& Spatial Communication}
\fancyfoot[C]{\small\thepage}

% ---------- Caption Style ----------
\captionsetup{
  font={small, sf},
  labelfont={bf},
  labelsep=period,
  skip=8pt
}

% ---------- Column types ----------
\newcolumntype{C}[1]{>{\centering\arraybackslash}p{#1}}
\newcolumntype{L}[1]{>{\raggedright\arraybackslash}p{#1}}

% ---------- Document ----------
\begin{document}

% Title block
\begin{center}
  {\Large\bfseries Comparison of Related Work}\\[4pt]
  {\large WebRTC, P2P Architectures, and Spatial Communication Systems}\\[6pt]
  \rule{0.6\textwidth}{1pt}
\end{center}

\vspace{1em}

\begin{table}[htbp]
\centering
\small
\setlength{\tabcolsep}{6pt}
\renewcommand{\arraystretch}{1.5}

\begin{tabularx}{\textwidth}{|C{0.6cm}|L{3.2cm}|L{3.3cm}|L{3.0cm}|L{3.3cm}|}
\hline

\textbf{No.} &
\textbf{Title, Year \& Venue} &
\textbf{Methodology} &
\textbf{Use Cases} &
\textbf{Research Gap} \\
\hline

1 &
Vivek Jain (2025), \textit{Server-Side Rendering vs Client-Side Rendering: A Comprehensive Analysis} &
Comparative analysis of SSR and CSR; performance evaluation and architectural differences &
Web application performance assessment; case studies of real rendering strategies &
Does not evaluate environmental impact or server load in distributed real-time spatial systems (focuses only on web pages) \\
\hline

2 &
T. Simao (2025), \textit{Evaluating Server-Side and Client-Side Rendering Performance} &
Empirical performance benchmarking under varying conditions (LCP, load metrics) &
SSR vs CSR performance across varied network and client conditions &
Focuses on web performance metrics; does not explore cross-domain distributed rendering or spatial engines \\
\hline

3 &
H. Zhang et al. (2024), \textit{Distributed Realtime Rendering in Decentralized Networks} &
Proposes decentralized rendering computation, task scheduling, and computation offloading &
Mobile Web AR distributed rendering; device-to-device collaborative computation &
Focuses on 3D AR object rendering; does not address 2D spatial communication or proximity-based client networks \\
\hline

4 &
Herzberger et al. (2023), \textit{Residency Octree: A Hybrid Approach for Scalable Web-Based Multi-Volume Rendering} &
Hybrid octree data structure for scalable client-side volume rendering &
Web-based multi-volume visualization with hierarchical cache and resolution control &
Does not address real-time collaboration or server-client partitioning for spatial interaction systems \\
\hline

5 &
S. Pati \& Y. Zaki (2025), \textit{Evaluating the Efficacy of Next.js — SSR vs CSR on Performance \& Accessibility} &
Comparative experimental study on framework strategies (SSR vs CSR) &
Web application performance metrics (FCP, TTI, user experience) across different network conditions &
Focuses on specific frameworks; does not evaluate distributed real-time world state management or spatial audio/video systems \\
\hline

\end{tabularx}

\caption{Comparison of Related Work in Rendering Strategies and Distributed Architectures}
\label{tab:related_work}
\end{table}

\end{document}